\providecommand{\CoHA}{\mathrm{CoHA}}
\providecommand{\Yp}{Y^{+}}
\providecommand{\dP}{\mathrm{dP}}
\providecommand{\Pic}{\mathrm{Pic}}
\providecommand{\Rep}{\mathrm{Rep}}
\providecommand{\Jac}{\mathrm{Jac}}
\providecommand{\Tot}{\operatorname{Tot}}
\providecommand{\Spec}{\operatorname{Spec}}
\providecommand{\Tr}{\operatorname{Tr}}
\providecommand{\Ext}{\operatorname{Ext}}
\providecommand{\kch}{\kappa_{\mathrm{ch}}}
\providecommand{\kcat}{\kappa_{\mathrm{cat}}}
\providecommand{\kBKM}{\kappa_{\mathrm{BKM}}}
\providecommand{\TP}{\mathcal{T}}
\providecommand{\C}{\mathbb{C}}
\providecommand{\Z}{\mathbb{Z}}
\providecommand{\P}{\mathbb{P}}
\providecommand{\Gm}{\mathbb{G}_m}

\section*{Local del Pezzo positive halves}

Let
\[
  S_n=\dP_n=\mathrm{Bl}_n(\P^2,p_1,\ldots,p_n),
  \qquad
  X_n=\Tot(K_{S_n}),
  \qquad 0\leq n\leq 8,
\]
with the points in general position. Then \(K_{S_n}^2=9-n\), and
\[
  K_{X_n}\simeq
  \pi^*(K_{S_n}\otimes K_{S_n}^{-1})
  \simeq \mathcal O_{X_n}.
\]
Thus \(X_n\) is a noncompact local Calabi-Yau threefold with compact
core \(S_n\).

A positive half attached to \(X_n\) is determined by the datum
\[
  (Q_n,W_n,\theta,o_n),
\]
where \((Q_n,W_n)\) is a \(3\)-Calabi-Yau quiver with potential,
\(\theta\) is a Bridgeland or King stability chamber, and \(o_n\) is
orientation data for the vanishing-cycle sheaf. For a dimension vector
\(\gamma\in\Z_{\geq0}^{Q_n^0}\), put
\[
  G_\gamma=\prod_{i\in Q_n^0}\mathrm{GL}_{\gamma_i}.
\]
The critical cohomological Hall algebra is
\[
  \Yp(X_n,\theta)
  :=
  \CoHA(Q_n,W_n)_\theta
  =
  \bigoplus_{\gamma\in\Z_{\geq0}^{Q_n^0}}
  H^{\mathrm{BM}}_{G_\gamma\times \Gm}\!
  \left(
    \Rep_\gamma(Q_n),
    \varphi_{\Tr(W_n)}\mathbb Q_{o_n}
  \right).
\]
Kontsevich-Soibelman construct the Hall product. Davison-Meinhardt give
the BPS integrality and PBW factorisation under the usual critical CoHA
orientation and finiteness hypotheses. The object above is an
associative \(E_1\) Hall algebra. In Calabi-Yau dimension \(3\), the
Hall comparison belongs to the \(E_1\)-chiral level of \(\Phi_3\). An
\(E_2\) structure enters only after adding a centre, a Drinfeld double
with a Hopf pairing, or a specified module category.

\subsection*{The local surface scalar}

The local scalar is a compact-core invariant. It is fixed by the
rank-one G\"ottsche-MNOP degree-zero normalization
\[
  \kch^{\mathrm{loc}}\!\left(\Tot(K_S)\right)
  =
  \kch(A_{\mathrm{loc}\,S})
  =
  \frac{1}{2}\chi_{\mathrm{top}}(S)
\]
for \(X=\Tot(K_S)\) with \(S\) smooth projective. For
\(S_n=\mathrm{Bl}_n(\P^2)\), the blow-up formula gives
\(\chi_{\mathrm{top}}(S_n)=3+n\), hence
\[
\begin{array}{c|ccccccccc}
n&0&1&2&3&4&5&6&7&8\\ \hline
K_{S_n}^{2}&9&8&7&6&5&4&3&2&1\\
\chi_{\mathrm{top}}(S_n)&3&4&5&6&7&8&9&10&11\\
\kch^{\mathrm{loc}}(X_n)&\frac32&2&\frac52&3&\frac72&4&\frac92&5&\frac{11}{2}
\end{array}
\]
For \(S_0=\P^2\),
\[
  \kch^{\mathrm{loc}}(K_{\P^2})
  =
  \frac{1}{2}\chi_{\mathrm{top}}(\P^2)
  =
  \frac{3}{2}.
\]
This is the chain-level local \(\P^2\) value recorded in
Theorem~\texttt{thm:local-p2-shadow} of the toric CY3 CoHA chapter and
in the Vol III stratification table. The calculation assigns the
local-surface \(\kch\) value coming from the compact core \(S_n\).
Compact Hodge supertraces, \(\kcat=\chi(\mathcal O_X)\), and
Borcherds-weight data are separate lanes. A \(\kBKM\) value requires a
Borcherds denominator with a Lorentzian lattice input; the local del
Pezzo surface supplies G\"ottsche-MNOP data.

The resolved conifold
\[
  \Tot(\mathcal O_{\P^1}(-1)\oplus \mathcal O_{\P^1}(-1))
  \longrightarrow \P^1
\]
lies in the curve-base local CY3 class, outside the local-surface class
\(\Tot(K_S)\to S\). Its direct McKay value is
\[
  \kch(X_{\mathrm{con}})=1.
\]

\subsection*{The \(dP_0\) quiver}

For \(S_0=\P^2\), Beilinson's exceptional collection
\[
  \langle \mathcal O,\mathcal O(1),\mathcal O(2)\rangle
\]
gives the standard \(dP_0\) tilting algebra. The total space
\(K_{\P^2}\) is the crepant resolution of
\(\C^3/(\Z/3)\), where the generator acts by
\[
  \zeta=\mathrm{diag}(\omega,\omega,\omega),
  \qquad
  \omega=e^{2\pi i/3}.
\]
Bridgeland-King-Reid, Theorem~1.2, identifies the crepant resolution
with the derived McKay category. The \(3\)-Calabi-Yau completion of the
Beilinson algebra is the quiver with vertices \(\Z/3\), arrows
\[
  a_i,b_i,c_i\colon i\longrightarrow i+1,
  \qquad i\in\Z/3,
\]
and cubic potential
\[
  W_0
  =
  \sum_{i\in\Z/3}
  a_i b_{i+1} c_{i+2}
  -
  a_i c_{i+1} b_{i+2}.
\]
Thus
\[
  \Yp(K_{\P^2},\theta)
  =
  \CoHA(Q_0,W_0)_\theta
\]
is the Hall positive half of the \(dP_0\) quiver Yangian in the
Rapcak-Soibelman-Yang-Zhao toric comparison. It is the local
\(\P^2\) analogue of the Schiffmann-Vasserot identification
\(\CoHA(\C^3)=Y^+(\widehat{\mathfrak{gl}}_1)\). The notation
\(\Yp\) here always means the positive Hall half. The full Yangian,
vertex algebra, and \(\mathcal W\)-algebra require the negative half,
Cartan factor, Hopf pairing, completion, centre, and vacuum-module
evaluation.

\subsection*{Toric \(dP_1,dP_2,dP_3\)}

For \(n=1,2,3\), \(S_n\) is toric and a toric brane-tiling or dimer
model supplies a quiver with potential \((Q_n,W_n)\). The number of
vertices is fixed by the rank of the even cohomology:
\[
  |Q_n^0|=\mathrm{rk}\,K_0(S_n)=\mathrm{rk}\,H^{\mathrm{ev}}(S_n,\Z)
  =n+3.
\]
A geometric blow-up adds the exceptional divisor class and adds one
vertex to the tilting presentation. Keller-Yang and
Derksen-Weyman-Zelevinsky mutation acts inside a fixed vertex set,
relating different toric chambers of the same \(X_n\). The change in
\(n\) is the geometric blow-up operation.

For each toric \(n\leq3\),
\[
  \Yp(X_n,\theta)=\CoHA(Q_n,W_n)_\theta
\]
is the critical CoHA of the chosen toric chamber. For two chambers related
by a Keller-Yang mutation at a loop-free, two-cycle-free
\(2\)-cycles, Keller-Yang Theorem~3.2 gives an equivalence of Ginzburg
dg categories. Transport of the CoHA additionally requires compatible
vanishing-cycle sheaves, orientation data, and stability chambers.
Iqbal-Kashani-Poor compute the refined
topological-vertex partition functions for these toric local del Pezzo
threefolds; after the standard DT/GV variable normalization, those
partition functions are the numerical chamber characters of the BPS
cohomology. A cohomological Hilbert-series equality requires the usual
purity and orientation-data hypotheses.

\subsection*{Non-toric \(dP_n\)}

For \(n\geq4\), the general-position blow-up of \(\P^2\) is non-toric.
A quiver model is supplied by a tilting bundle or by an explicit
critical atlas. If \(\TP_n\) is a tilting bundle on \(S_n\) and
\(\pi:X_n\to S_n\) is the projection, then
\[
  \mathrm{End}_{X_n}(\pi^*\TP_n)
  =
  \bigoplus_{m\geq0}
  \mathrm{Hom}_{S_n}\!\left(\TP_n,\TP_n\otimes K_{S_n}^{-m}\right)
\]
is the rolled-up helix algebra. The Hall positive half is formed from a
Jacobi presentation
\[
  \Jac(Q_n,W_n)
\]
or from the corresponding Ginzburg dg algebra \(\Gamma(Q_n,W_n)\) as
the \(3\)-Calabi-Yau presentation of the local surface category.

The root lattice controlling the del Pezzo chamber symmetry is
\[
  R_n=\{\alpha\in\Pic(S_n):\alpha^2=-2,\ \alpha\cdot K_{S_n}=0\}.
\]
For \(3\leq n\leq8\) one has
\[
\begin{array}{c|cccccc}
n&3&4&5&6&7&8\\ \hline
R_n&A_1\oplus A_2&A_4&D_5&E_6&E_7&E_8 .
\end{array}
\]
The Weyl group \(W(R_n)\) acts on \(\Pic(S_n)\) and on the set of
\((-1)\)-curves. Where the categorical representatives exist, the
corresponding mutation braid group acts on exceptional collections.
This supplies chamber symmetry for the BPS quivers. Recognition of
\(\Yp(X_n,\theta)\) as a positive half of an affine \(E_n\)-type
Yangian for \(n=6,7,8\) requires the following inputs:
\[
  (Q_n,W_n),\quad
  \text{orientation data},\quad
  \text{a PBW theorem},\quad
  \text{Cartan/current normalization},\quad
  \text{character comparison}.
\]
The \(E_6,E_7,E_8\) symmetry visible in the del Pezzo lattice and in
the Iqbal-Kashani-Poor character calculations gives character-level
support for that recognition.

\subsection*{BPS character}

For a constructed datum \((Q_n,W_n,\theta,o_n)\), the
Davison-Meinhardt integrality theorem gives a PBW-type BPS
factorisation of the completed critical CoHA:
\[
  \operatorname{gr}\CoHA(Q_n,W_n)_\theta^{\wedge}
  \cong
  \mathrm{Sym}\!\left(
    \bigoplus_{\gamma}
    \mathrm{BPS}_{\theta,\gamma}\otimes H^\bullet(B\Gm)
  \right)
\]
after ordering by the central charge. Its Euler character is
\[
  \chi_{\mathrm{CoHA}}(q,Q)
  =
  \sum_{\gamma}
  \chi\!\left(
    H^{\mathrm{BM}}_{G_\gamma\times\Gm}
    (\Rep_\gamma(Q_n),\varphi_{\Tr(W_n)}\mathbb Q_{o_n})
  \right) q^{\gamma_0}Q^\gamma .
\]
In the toric chambers \(n=0,1,2,3\), virtual localization and the
refined topological vertex compute the same DT/GV chamber character
\[
  \chi_{\mathrm{CoHA}}(q,Q)
  =
  Z_{\mathrm{top}}^{X_n,\theta}(q,Q)
\]
after the chosen DT/GV normalization. For \(n\geq4\), the same equality
requires the quiver-with-potential model, compact-support orientation
data, and a wall-crossing comparison with the chosen BPS chamber.

\subsection*{Proof boundary}

Outside the toric range, the remaining inputs are concrete: explicit
\((Q_n,W_n)\) presentations for the selected tilting models, orientation
data for compact-support vanishing-cycle cohomology, purity or a
replacement for Hilbert-series lifts from Euler characters, and a
chamber-by-chamber comparison with the refined topological-string
amplitudes. Passage from a positive half to a full Yangian or
\(\mathcal W\)-object adds the negative half, Cartan part, Hopf pairing,
completion, Drinfeld double or centre, and a specified representation.
